\section{Esperança de Variáveis Discretas}

A esperança de uma variável aleatória é definida matematicamente como a média ponderada
dos valores assumidos por essa variável. Uma interpretação para esperança é que ela é
o valor médio de todas as realizações da variável ao se repetir o experimento aleatório
um número infinito de vezes. Essa interpretação será justificada posteriormente com
a \textbf{Lei dos Grandes Números}.

\begin{def:esperanca discreta}
Seja $X$ uma v.a. discreta com f.m.p. $p_X$. A esperança $E[X]$ de $X$ é definida tal como
\[
	E[X] = \sum_{x \in \Omega_X} x \cdot p_X(x)
\]
\end{def:esperanca discreta}

A partir do conhecimento da distribuição dos valores de uma variável, é possível ter
conhecimento da distribuição dos valores de uma função aplicada sobre essa variável.
Se $X: \Omega \to \mathbb{R}$ e $f: \mathbb{R} \to \mathbb{R}$, então $f(X) = f \circ X$
é uma variável aleatória por definição. Digamos que $f(X) = Y$, então a massa de
probabilidade $Y = y$ será a soma das probabilidades de $X = x$ tal que $x \in \Omega_X$
tem $y$ como imagem em $f$, já que as realizações de uma v.a. são sempre disjuntas.
\[
	P(Y = y) = \sum_{\substack{x \in \Omega_X \\ f(x) = y}} P(X = x)
\]
\noindent
Tão logo conhecendo a dsitribuição de $f(X)$ a partir da distribuição de $X$, então poderemos
expressar a esperança de $f(X)$ em termos da própria f.m.p. de $X$.

\begin{prop:esperanca de funcao}
Seja $X$ uma v.a. discreta com f.m.p. $p_X$, e $f: \mathbb{R} \to \mathbb{R}$
uma função. A esperança da variável aleatória definida por $f(X) = f \circ X$ é
dada por
\[
	E[f(X)] = \sum_{x \in \Omega_X} f(x) \cdot p_X(x)
\]
\begin{proof}
	Seja $\Omega_f$ o suporte de $f(X)$ e $p_f$ a sua f.m.p.. Temos que
	\begin{align*}
		\sum_{x \in \Omega_X} f(x) \cdot p_X(x) & = \sum_{y \in \Omega_f} \sum_{\substack{x \in \Omega_x   \\ f(x) = y}} y \cdot p_X(x)\\
		                                        & = \sum_{y \in \Omega_f} y \sum_{\substack{x \in \Omega_x \\ f(x) = y}} p_X(x)
	\end{align*}
	Fixando $y$, observamos que
	\[
		p_f(y) = \sum_{\substack{x \in \Omega_x \\ f(x) = y}} p_X(x).
	\]
	Logo
	\[
		E[f(X)] = \sum_{y \in \Omega_f} y \cdot p_f(y).
	\]
\end{proof}
\end{prop:esperanca de funcao}

A esperança satisfaz as propriedades de linearidade. Isso significa que, se uma v.a.
é definida como uma constante, como produto de uma v.a. por uma constante ou como a soma de
outras duas v.a.'s, podemos facilmente obter a esperança da variável derivada a partir das
variáveis que a compõe. Ressaltemos que constantes podem ser vistas como v.a.'s com
suporte unitário.

\begin{prop:linearidade da esperanca discreta}
Sejam $X$ e $Y$ variáveis aleatórias. A esperança satisfaz:
\begin{enumerate}
	\item se $X$ tem suporte unitário igual a $\{c\}$, então $E[X] = c$.
	\item $E[\lambda X] = \lambda E[X]$.
	\item $E[X + Y] = E[X] + E[Y]$.
\end{enumerate}
\begin{proof}
	Para (1), basta observar que $E[X] = c \cdot 1 = c$.

	Para (2),  usemos a definição
	de esperança:
	\[
		E[\lambda X] = \sum_{x \in \Omega_X} \lambda x \cdot  p_X(x)
		= \lambda \sum_{x \in \Omega_X} x \cdot p_X(x)
		= \lambda E[X].
	\]
	Para provar (3), temos que
	\[
		E[X + Y] = \sum_{x \in \Omega_X} \sum_{y \in \Omega_Y} (x+y) P(X = x, Y = y),
	\]
	e expandindo,
	\begin{equation}
		\label{eq:dem esperanca linear}
		E[X + Y] =
		\sum_{x \in \Omega_X} x \sum_{y \in \Omega_Y} P(X = x, Y = y) +
		\sum_{y \in \Omega_Y} y \sum_{x \in \Omega_X} P(X = x, Y = y) .
	\end{equation}
	Já que
	\[
		\sum_{y \in \Omega_Y} P(X = x, Y = y) = \sum_{y \in \Omega_Y} P(X = x | Y = y) P(Y = y),
	\]
	então pelo Teorema da Probabilidade Total,
	\[
		\sum_{y \in \Omega_Y} P(X = x, Y = y) = P(X = x)
	\]
	e, de forma análoga,
	\[
		\sum_{x \in \Omega_X} P(X = x, Y = y) = P(Y = y),
	\]
	então podemos simplificar a expressão \ref{eq:dem esperanca linear} para
	\[
		E[X + Y] = \sum_{x \in \Omega_X} x \cdot P(X = x) + \sum_{y \in \Omega_Y} y \cdot P(Y=y),
	\]
	e, finalmente
	\[
		E[X + Y] = E[X] + E[Y]
	\]
\end{proof}
\end{prop:linearidade da esperanca discreta}
